% !TeX spellcheck = en
\documentclass[12pt, a4paper]{extarticle}

\input{/Users/arhammahajan/Developer/dotfiles/latex-template/preamble.tex}
\SetHeader{Computer Aided Design and Manufacturing}{Unit-1}

\date{\today}

\begin{document}

\maketitle
\tableofcontents
\newpage

% --- Start writing your notes below ---

\section{Introduction}
[Introduction text goes here...]

\section{[Main Topic 1]}
[Text content...]

\subsection{[Subtopic 1.1]}
[Text content...]

% Example of the Gray Secondary Box for Equations
\begin{theorem}{Important Equation}{}
$$
E = mc^2
$$
\end{theorem}

\begin{theoremc}
	This is a continuation of the above theorem
\end{theoremc}

\begin{corollary}{Edge cases}{}
	This is a corollary to the above
\end{corollary}

\begin{itemize}
    \item Bullet point 1
    \item Bullet point 2
\end{itemize}

\subsection{[Subtopic 1.2]}
[Text content...]

\begin{defn}{Special theory of relativity}{}
	This is an example definition
\end{defn}

\begin{prop}{NVIDIA is better than AMD}{}
	This is an example proposition
\end{prop}

\begin{example}
\[
\det\begin{pmatrix} 2 & 3 \\ 6 & 10 \end{pmatrix} = 2(10) - 3(6) = 2
\]
\end{example}

\begin{property}{Properties of Transpose}{}
\begin{enumerate}
	\item $(A^T)^T = A$
	\item $(A + B)^T = A^T + B^T$
	\item $(cA)^T = cA^T$
	\item $(AB)^T = B^TA^T$ (note the reversal of order)
\end{enumerate}
\end{property}


% Example of Multiple images in one row
% \begin{figure}[H]
%     \centering
%     \begin{minipage}{0.45\textwidth}
%         \centering
%         \includegraphics[width=\linewidth]{path/to/image-1.png}
%         \caption{Caption-1}
%     \end{minipage}
%     \hfill
%     \begin{minipage}{0.45\textwidth}
%         \centering
%         \includegraphics[width=\linewidth]{path/to/image-2.png}
%         \caption{Caption-2}
%     \end{minipage}
% \end{figure}

\begin{claim}{I am not claiming anything}{}
	Example
\end{claim}

\begin{remark}[Non-commutativity]
	In general, matrix multiplication is \textbf{not commutative}: $AB \neq BA$.
\end{remark}

\begin{qs}{Very important question}{}
	What is the mass of the earth?
\end{qs}

The solution to the above is...

\begin{note}
	Keep this in mind
\end{note}

\end{document}